\begin{tablalafrox}{Matriz de comunicaciones del proyecto. Fuente: elaboración a partir de (Distribuidora Puelche S.A., 2026a, Bases Administrativas, art. 71, pp. 37--38), (Subdocumento 6, Anexo 6.A) y (Formulario T-14, paquetes 1.2.5, 1.3.2, 1.7.2 a 1.7.7, 9.5.2 y 9.5.3).}{tab:tabla-comunicaciones}%
  {F{0.20} F{0.17} F{0.20} F{0.17} Y}%
  {\cab{Comunicación} & \cab{Responsable} & \cab{Destinatarios} & \cab{Frecuencia o momento} & \cab{Canal y registro}}%
  [Matriz de comunicaciones del proyecto]
  Seguimiento de actividades, compromisos y bloqueos & Jefe de Proyecto & Equipos de ambas partes & Semanal (1.2.5) & Reunión presencial o remota; informe semanal en el espacio colaborativo \\
  Avance de entregables, desviaciones y riesgos & Jefe de Proyecto & Comité de Proyecto & Quincenal (1.7.3) & Sesión con acta en dos días hábiles (RT-19.09) \\
  Informe de avance con valor ganado & Jefe de Proyecto & Comité Ejecutivo y Contraparte Técnica & Mensual (1.7.6) & Informe con SPI y CPI; tablero del proyecto (1.7.7) \\
  Decisiones y riesgos de arquitectura & Arquitecto de Solución & Comité de Arquitectura & Mensual (1.7.4) & Sesión con acta; registro de decisiones del SD4 \\
  Niveles de servicio, incidentes y capacidad & Líder de Operación / SRE & Comité de Operación & Mensual desde el mes 13 (1.7.5) & Sesión con acta; informe de niveles de servicio \\
  Solicitud o decisión requerida del CLIENTE & Responsable de la actividad & Contraparte Técnica & Diez días hábiles antes de la actividad, con respuesta a cinco & Solicitud registrada con ID, evidencia y fecha límite (T-9, sección 1.6) \\
  Cambios de la próxima ola: qué cambia, cuándo y qué se espera de cada persona & Líder de Implantación & Personal de bodega, preventistas y tripulaciones & Dos semanas y dos días antes de cada ola & Reunión de inicio de jornada con el jefe de frente; afiche en andén y bodega; aviso en el terminal \\
  Turno de noche de bodega & Jefe de frente de bodega & Preparadores, jefa de bodega y tutor de turno & Cada cambio de turno durante olas y estabilización; quincenal en régimen (9.5.2) & Reunión de cambio de turno a las 22:00; registro de asistencia; guía rápida en el puesto \\
  Preventistas y tripulaciones en ruta & Jefe de frente de ruta & Preventistas, conductores y peonetas & Inicio de jornada en la semana de la ola; semanal en marcha blanca & Reunión breve en el andén; aviso en el terminal que se lee sin señal \\
  Conductores de transportistas & Jefe de frente de ruta & Conductores externos y su empresa & Al asignarse a un viaje (9.5.3) & Instrucción en el andén y guía impresa; aviso escrito a la empresa transportista \\
  Planificador de rutas & Líder de Implantación & Planificador y su ayudante & Talleres en los meses 1 a 3; quincenal durante la construcción de M4 & Sesión individual con acta firmada de las reglas validadas (2.1.4) \\
  Sindicato de choferes & Jefe de Proyecto & Directiva sindical & Antes de usar GPS para control de jornada y ante cambios de finalidad o privacidad que afecten a conductores & Mesa de trabajo con acta (9.7.2) \\
  Cadenas del canal moderno & Arquitecto de Solución & Responsable técnico de cada cadena & Acuerdo durante el levantamiento (2.1.5); seguimiento quincenal desde el inicio de la certificación & Mesa técnica por cadena; registro del perfil, acceso, agenda y pruebas (5.3.1 y 5.3.2) \\
  Clientes, proveedores y autoridad sanitaria & Líder de Implantación & Clientes por canal, 180 proveedores y autoridad sanitaria & Un mes antes del cambio que les afecta; a solicitud en el caso de la autoridad & Preventista o ejecutivo comercial; comunicado formal; canal de la Jefa de Calidad \\
\end{tablalafrox}
